\begin{abstract}
We simulate the silicon-nitride (SiN) photonic circuits that are to steer the transport tweezers (1064.4~nm) and
switch the addressing beams (894.6~nm) of a Cs neutral-atom processor, and the system built from them, up to
arrays of $10^6$ sites. The 400~nm SiN platform meets the proposal's component targets; two targets need
restating: the steering update rate should be $\geq 200$~kHz for moves at 1.2~mK, and position noise should be
specified as a spectral density at the trap frequency. Independent motion of about 20 tweezers along a row
requires a sliding-window focal-plane phased array. Its optimised design (128 ports, 192 emitters at 1.5~\textmu m
pitch) needs about 57~mW per 1.2~mK tweezer once facet loss and a glued collimator are included. Chips stacked in a
16-row anamorphic relay give independent rows on a 4~\textmu m lattice without per-chip frequency shifters, using
common phase noise on the seed laser and fibre delays. In three-dimensional transport with handoff, 3,000 of 3,000
simulated atoms are delivered along lanes between the rows, with 0.12~\textmu K of heating per handoff (about 0.3~\textmu K per move) at a
moving depth of 0.3~mK when the ramps last 200~\textmu s. At $10^6$ sites, a block of 128 engines holds a
time-averaged defect fraction of about $2\times10^{-4}$, set by vacuum and imaging losses. Light that forms tweezers
away from the row centre detunes a refractive cylindrical collimator, which then serves only the central half
($\pm45$~\textmu m) of a 180~\textmu m segment. Two mirrors in a glued glass block, designed here, collimate exactly
at every such angle and restore the full segment (peak $\geq0.98$ to the slit edge), at the cost of a 10~nm rms
mirror form tolerance and a coating-dependent circular polarisation ($|V|\leq0.04$ at 90~\textmu m). All values are simulation
outputs or derivations within stated models and should be reviewed before use.
\end{abstract}
